
\section{引言}
\subsection{问题背景}

导电复合材料在新能源材料、半导体器件和电磁屏蔽等领域具有广泛的应用前景。此类材料由绝缘基体和分散在其中的导电填料组成，其导电性能取决于填料颗粒在基体中能否形成贯穿的导电通路。理解填料几何形状、填充比例与导电概率之间的定量关系，对于材料设计和工艺优化具有重要的理论指导意义。

本研究将实际导电复合材料简化为一个理想化的模型系统——微构体。微构体是边长为10000nm的立方体区域，内部默认填充绝缘介质。导电填料分为两类：金属A为长5000nm、底面半径30nm的直圆柱体，具有极高的长径比（约83:1）；金属B为半径200nm的球体，长径比为1:1。导电机制基于量子隧穿效应：当两颗粒表面之间的最短距离不超过1.8nm时，认为两者之间形成导电连接。微构体的左右两侧面（X=-5000和X=+5000）分别作为电极面，若存在从一侧电极到另一侧电极的连通路径，则判定微构体导通。

该模型系统涵盖了导电复合材料研究中的几个核心科学问题：填料几何形状对渗透阈值的影响、填充比例与导通概率的定量关系、以及双填料体系的成本优化。通过建立数学模型和数值模拟方法，可以系统揭示这些关系。

\subsection{研究意义}

本研究的意义体现在以下方面。在方法论层面，本文建立了从三维几何建模到蒙特卡洛模拟再到优化决策的完整计算框架，为导电复合材料的设计提供了定量分析工具。在物理认知层面，本文揭示了高长径比圆柱形填料具有超低渗透阈值的特性——体积分数仅为0.0551\%即可实现90\%以上的导通概率，这一发现对低填料含量导电复合材料的开发具有指导价值。在工程应用层面，本文的双金属成本优化分析为材料配方设计提供了经济性视角的决策参考。

\subsection{问题重述}

\textbf{问题一：}微构体中导电介质仅含金属A。若仅有一个金属A圆柱体，给定其在微构体中的两端点坐标，分析该微构体是否导通，并计算导通概率。

\textbf{问题二：}微构体中仅含金属A圆柱体，金属A的体积分数为0.50\%。在该条件下计算微构体的导通概率，并进一步计算体积分数为0.60\%、0.70\%和1.00\%时的导通概率。

\textbf{问题三：}微构体中仅含金属A圆柱体，在导通概率不低于90\%的前提下，求金属A的最小体积分数。

\textbf{问题四：}微构体中同时含金属A圆柱体和金属B球体。金属A成本为1.05元/$\mu$m$^3$，金属B成本为0.05元/$\mu$m$^3$。在导通概率不低于90\%的前提下，求使总成本最低的两种金属的最优数量。
